\begin{textsample}{POS Dim 2 – persona\_gemini – Score 19.00 – t221\_gemini.txt}  \label{ex:f2_pos_048}
Awa Traore reflects on COP26 with mixed feelings, pointing to the profit-driven \textbf{corporations} and \textbf{leaders} who are destroying \textbf{ecosystems} and threatening Indigenous \textbf{lands}.
She calls on the Global North to acknowledge the role of exploitative capitalism in the \textbf{climate} \textbf{crisis}, while highlighting that the Global South prioritizes \textbf{people} over carbon offsets.
COP26 failed on its key objectives.
It did not secure a \textbf{path} to \textbf{limit} warming to 1.5C, failed to deliver the promised \$100 billion in \textbf{climate} \textbf{finance}, and did not establish a \textbf{fund} for \textbf{loss} and damage.
Instead, the \textbf{summit} promoted unfair carbon trading \textbf{systems} that harm \textbf{communities} in the Global South.
Despite these \textbf{outcomes}, there is hope to be found in \textbf{people} \textbf{power}.
Indigenous activists and young \textbf{people} are confronting \textbf{leaders}, \textbf{demanding} not just \textbf{climate} action, but economic and racial \textbf{justice} as well.
In looking ahead, Awa Traore calls for stronger \textbf{resistance} to \textbf{center} \textbf{people} and the \textbf{planet} at COP27.

% matched lemmas: center, climate, community, corporation, crisis, demand, ecosystem, finance, fund, justice, land, leader, limit, loss, outcome, path, people, planet, power, resistance, summit, system
\end{textsample}
